\documentclass[11pt, a4paper]{article}

\usepackage[margin=2.5cm]{geometry}
\usepackage{amsmath, amsthm, amssymb}
\usepackage{mathtools}
\usepackage{hyperref}
\usepackage{cleveref}
\usepackage{booktabs}
\usepackage{microtype}
\usepackage{listings}
\usepackage{xcolor}

% ── Theorem environments ───────────────────────────────────────────────────
\newtheorem{theorem}{Theorem}[section]
\newtheorem{lemma}[theorem]{Lemma}
\newtheorem{proposition}[theorem]{Proposition}
\newtheorem{corollary}[theorem]{Corollary}
\newtheorem{definition}[theorem]{Definition}
\newtheorem{remark}[theorem]{Remark}

\theoremstyle{remark}
\newtheorem*{note}{Note}

% ── Macros ─────────────────────────────────────────────────────────────────
\newcommand{\R}{\mathbb{R}}
\newcommand{\N}{\mathbb{N}}
\newcommand{\Z}{\mathbb{Z}}
\newcommand{\wmax}{W_{\max}}
\newcommand{\crowd}{\mathrm{crowd}}
\newcommand{\pen}{\mathrm{pen}}
\newcommand{\CVaR}{\mathrm{CVaR}}
\newcommand{\Var}{\mathrm{Var}}
\newcommand{\E}{\mathbb{E}}

\title{%
    \textbf{Namma Metro Time-Dependent Routing Engine}\\[0.5em]
    \large Formal Correctness Proofs and Mathematical Analysis
}
\author{Darshan Rajagoli}
\date{\today}

% ──────────────────────────────────────────────────────────────────────────
\begin{document}
\maketitle
\tableofcontents
\newpage

% ══════════════════════════════════════════════════════════════════════════
\section{Formal Graph Model}
% ══════════════════════════════════════════════════════════════════════════

\begin{definition}[Transit Graph]
    Let $G = (V, E)$ be a directed graph where:
    \begin{itemize}
        \item $V = \{v_1, v_2, \ldots, v_n\}$ is the set of transit stops
              (stations on the Namma Metro Purple and Green Lines), $|V| = n$.
        \item $E \subseteq V \times V \times \N$ is the set of time-stamped
              directed service legs. Each edge $e = (u, v, t_d) \in E$ encodes
              a departure from stop $u$ at absolute time $t_d$
              (seconds past midnight), arriving at stop $v$ at time
              $t_a = t_d + \tau(u, v)$, where $\tau: V \times V \to \N^+$
              is the travel duration.
    \end{itemize}
\end{definition}

\begin{definition}[Time-Dependent Edge Weight]
    The time-dependent weight function $w: V \times V \times \N \to \N$ is
    defined as
    \[
        w(u, v, t) \;=\; \min \bigl\{ t_a - t \;\mid\; (u, v, t_d) \in E,\; t_d \ge t \bigr\},
    \]
    representing the minimum elapsed time from $t$ until arrival at $v$ by
    boarding any service from $u$ that departs at or after $t$.
\end{definition}

\begin{definition}[Bi-Criteria Edge Weight]
    In the bi-criteria extension, each edge carries both a temporal cost and
    a crowd cost. The \textbf{full composite cost} used by
    \texttt{select\_optimal\_departure} is:
    \[
        c(e, t) \;=\; \tau(e) + \lambda \cdot \crowd(e) + \pen(u, v, t),
    \]
    where $\tau(e)$ is travel time, $\lambda \ge 0$ is the crowd-aversion
    parameter, $\crowd(e)$ is the static crowd weight of edge $e$, and
    $\pen: V \times V \times \N \to \R_{\ge 0}$ is a time-dependent
    surcharge modelling wait-time cost derived from IUDX AFC gate data.

    \begin{note}
    The $\lambda \cdot \crowd(e)$ term is \textbf{time-independent}: it is a
    static property of the edge, not a function of $t$. Therefore the FIFO
    analysis applies to the full composite:
    $\frac{d}{dt}[\lambda \cdot \crowd(e)] = 0$, and the derivative constraint
    reduces to a condition on $\pen$ alone.
    \end{note}
\end{definition}

% ══════════════════════════════════════════════════════════════════════════
\section{The FIFO Property and Penalty Constraints}
% ══════════════════════════════════════════════════════════════════════════

\begin{definition}[FIFO Property]
    A time-dependent edge $(u, v)$ satisfies the \emph{First-In-First-Out
    (FIFO) property} if and only if for all $t_1 \le t_2$:
    \[
        t_1 + w(u, v, t_1) \;\le\; t_2 + w(u, v, t_2).
    \]
    Equivalently, the arrival-time function $f(t) = t + w(u, v, t)$ is
    non-decreasing. A network is FIFO if all its edges satisfy this property.
\end{definition}

\begin{remark}
    FIFO is the necessary and sufficient condition for Dijkstra's algorithm to
    maintain its optimality subpath property on time-dependent graphs. Violation
    would allow a later departure to arrive earlier, invalidating the greedy
    relaxation step.
\end{remark}

\subsection{The FIFO Lemma for the Full Composite Cost}

\begin{lemma}[Penalty Derivative Constraint]
    \label{lem:fifo-pen}
    Consider the full composite arrival function for edge $(u,v)$ at query
    time $t$:
    \[
        f(t) \;=\; t + w(u, v, t) + \pen(u, v, t),
    \]
    where the $\lambda \cdot \crowd$ term is time-independent and drops out of
    the derivative. Then $f$ satisfies FIFO if the following \textbf{sufficient
    condition} holds:
    \[
        \pen(u, v, t_2) - \pen(u, v, t_1)
        \;\ge\; -(t_2 - t_1)
        \quad \text{for all } t_1 \le t_2.
    \]
    For differentiable $\pen$, the condition becomes $\frac{d}{dt}\pen \ge -1$.
    \begin{note}
    This is a \emph{sufficient} condition for FIFO, not a necessary one.
    The intra-interval necessity analysis (see proof) shows the tighter condition
    $\pen(t_2) - \pen(t_1) \ge 0$ within each inter-departure interval.
    For the scaffold where $\pen = 0$ identically, both conditions hold
    trivially. The sufficient condition $\ge -(t_2-t_1)$ is the operationally
    useful constraint because it is the one verified by \texttt{test\_fifo\_invariant.cpp}.
    \end{note}
\end{lemma}

\begin{proof}
    \textbf{Finite-difference formulation.}
    $w(u,v,t)$ is piecewise-constant between consecutive departure events and
    is not differentiable at departure points; applying $d/dt$ is formally
    invalid at those points. The correct statement uses finite differences.

    Between consecutive departures $t_d$ and $t_{d+1}$, the wait function
    $w(u,v,t) = t_d - t + \tau$ decreases at rate $-1$, so
    $w(u,v,t_2) - w(u,v,t_1) = -(t_2 - t_1)$ for $t_1, t_2$ in the same
    inter-departure interval. At a departure event $t_d$, $w$ jumps
    discontinuously upward (wait resets to the next service interval), so
    $w(u,v,t_2) - w(u,v,t_1) \ge -(t_2 - t_1)$ holds in all cases.

    \textbf{Sufficiency.}
    FIFO requires $f(t_2) - f(t_1) \ge 0$ for $t_1 \le t_2$:
    \begin{align*}
        f(t_2) - f(t_1)
        &= (t_2 - t_1) + [w(u,v,t_2) - w(u,v,t_1)]
           + [\pen(t_2) - \pen(t_1)] \\
        &\ge (t_2 - t_1) - (t_2 - t_1) + [\pen(t_2) - \pen(t_1)] \\
        &= \pen(t_2) - \pen(t_1).
    \end{align*}
    If $\pen(t_2) - \pen(t_1) \ge -(t_2 - t_1)$, then
    $f(t_2) - f(t_1) \ge 0$. \checkmark

    \textbf{Intra-interval analysis (tighter necessary condition).}
    Within a single inter-departure interval $(t_{d-1}, t_d)$, $w$ decreases
    at rate exactly $-1$: $w(u,v,t_2) - w(u,v,t_1) = -(t_2-t_1)$.
    Substituting into $f(t_2) - f(t_1) \ge 0$:
    \[
        0 \;\le\; (t_2-t_1) - (t_2-t_1) + [\pen(t_2)-\pen(t_1)]
          \;=\; \pen(t_2)-\pen(t_1),
    \]
    giving the intra-interval necessary condition
    $\pen(t_2) - \pen(t_1) \ge 0$ for $t_1 \le t_2$ within the interval.
    This is \emph{stronger} than $\ge -(t_2-t_1)$.
    For pairs straddling a departure event, the constraint involves the service
    headway $h = t_d - t_{d-1}$ rather than $t_2 - t_1$, and is weaker.
    The sufficient condition $\ge -(t_2-t_1)$ is the uniform constraint that
    covers all cases conservatively and is operationally checkable. \checkmark

    \textbf{Scope of the FIFO guarantee.}
    The FIFO condition ensures that $\phi(t) = t_{\text{dep}}(t) + \tau$
    is non-decreasing in query time $t$, where $t_{\text{dep}}(t)$ is the
    departure selected at query time $t$. This is what Dijkstra's optimality
    requires: a passenger who arrives at node $u$ later cannot be routed to
    node $v$ earlier, so the greedy relaxation order is valid.

    \textbf{Correctness of Pareto-Dijkstra (consistency under extension).}
    For the label-correcting bi-criteria algorithm, FIFO is necessary but not
    sufficient for correctness. The additional required property is
    \emph{consistency under extension}: if label $L_1 \preceq L_2$ at node $u$
    (i.e.\ $L_1$ weakly dominates $L_2$ in both dimensions), then for any edge
    $(u,v)$, the extensions $\mathrm{ext}(L_1) \preceq \mathrm{ext}(L_2)$ at
    $v$. For additive cost accumulation (arrival\_time and crowd\_cost both
    accumulate as sums along the path):
    \[
        (t_1 \le t_2 \;\wedge\; c_1 \le c_2)
        \;\Rightarrow\;
        (t_1 + \Delta t \le t_2 + \Delta t \;\wedge\; c_1 + \Delta c \le c_2 + \Delta c)
    \]
    for any $\Delta t, \Delta c \ge 0$. This holds trivially for non-negative
    edge weights, which is guaranteed by the minimum travel\_time $\ge 1$\,s
    constraint enforced in \texttt{graph\_builder.cpp}. Together, FIFO and
    consistency under extension establish the correctness of the label-correcting
    Pareto-Dijkstra for this routing engine.
\end{proof}

% ══════════════════════════════════════════════════════════════════════════
\section{Suboptimality of Strict Next-Train Boarding in Off-Peak Hours}
% ══════════════════════════════════════════════════════════════════════════

\begin{proposition}[Strict Next-Train Suboptimality]
    \label{prop:strict-suboptimal}
    Let $e_1 = (u, v, t_1)$ and $e_2 = (u, v, t_2)$ with $t_1 < t_2$ be two
    consecutive departures on the same service. Suppose the pen function
    exhibits an off-peak precipitous drop such that:
    \[
        \pen(u, v, t_2) < \pen(u, v, t_1) - (t_2 - t_1).
    \]
    Then the strict next-train policy (always boarding $e_1$) produces a
    composite cost strictly greater than boarding $e_2$:
    \[
        w'(u, v, t_1) \;>\; (t_2 - t_1) + w'(u, v, t_2).
    \]
\end{proposition}

\begin{proof}
    For consecutive departures on the same service, $\tau_1 = \tau_2 = \tau$
    and $\crowd_1 = \crowd_2$ (same service leg, consecutive timetable entries).
    The implementation composite (no wait time — see \texttt{routing.hpp} §3)
    is $c(e) = \tau + \lambda\cdot\crowd + \pen$.
    Under the equal-$\tau$, equal-$\crowd$ assumption:
    \[
        c(e_1) - c(e_2) = \pen_1 - \pen_2.
    \]
    By the hypothesis $\pen_1 - \pen_2 > t_2 - t_1 > 0$, so
    $c(e_1) > c(e_2)$. Therefore strict next-train boarding (selecting $e_1$)
    yields strictly higher composite cost than boarding $e_2$.

    \begin{note}
    Both $e_1$ and $e_2$ remain \emph{Pareto-optimal} in (arrival\_time,
    crowd\_cost) space: $e_1$ arrives earlier ($t_1 + \tau < t_2 + \tau$) with
    higher crowd, while $e_2$ arrives later with lower crowd. Neither dominates
    the other. The proposition establishes that strict next-train boarding is
    suboptimal for the \emph{scalar composite objective}, not that $e_1$ should
    be excluded from the Pareto frontier. The multi-label algorithm correctly
    retains both labels.
    \end{note}
\end{proof}

\begin{corollary}[Bounded-Wait Lookahead Resolves This]
    The Bounded-Wait Lookahead policy, which evaluates the next $k$ departures
    within the window $[t, t + \wmax]$ and selects $e^* = \arg\min_i w'(u,v,t_i)$,
    provably selects $e_2$ over $e_1$ in the scenario of
    \Cref{prop:strict-suboptimal}, provided $t_2 \le t + \wmax$.
\end{corollary}

% ══════════════════════════════════════════════════════════════════════════
\section{The Markowitz Analogy and Its Limitations}
% ══════════════════════════════════════════════════════════════════════════

\subsection{The Analogy}

The bi-criteria routing objective:
\[
    \min_{\text{path } P} \;\; \text{travel\_time}(P) + \lambda \cdot \crowd(P)
\]
is \textbf{structurally analogous} to the classical Markowitz mean-variance
efficient frontier~\cite{markowitz1952}. In both problems, varying $\lambda$
traces a frontier of achievable (cost, risk) tradeoffs:

\begin{note}
The analogy is informal. The transit objective is a \emph{deterministic linear
scalarization} over fixed edge properties; Markowitz minimizes a \emph{quadratic
form} $\mathbf{w}^\top \Sigma\, \mathbf{w} - \mu^\top\mathbf{w}/\lambda$ over
a distribution of random returns. The $\lambda$ parameter plays the role of
risk-aversion in both settings, but there is no probability distribution, no
variance, and no covariance matrix in this model. Writing the transit objective
as $\E[\text{travel\_time}] + \lambda \cdot \sigma[\crowd]$ is informal notation
for the deterministic scalarization; it does not imply a stochastic model.
The classical Markowitz problem is linear in $\sigma^2$ (variance, quadratic in
$\sigma$); this objective is linear in $\crowd$ directly. For a rigorous analogy
one would write: varying $\lambda$ traces the convex hull of the Pareto frontier,
analogous to how varying risk-aversion traces the efficient frontier. Both
characterize the same tradeoff geometry without requiring structural isomorphism.
\end{note}

\begin{remark}[Lambda-Scalarisation Recovers Only the Convex Hull]
    \label{rem:convex-hull}
    A critical limitation of the Markowitz analogy is that varying $\lambda$
    over $[0, \infty)$ traces only the \textbf{convex hull} of the Pareto
    frontier, not the full frontier. For a discrete graph, the Pareto frontier
    is a finite set of points in $(\text{time}, \text{crowd})$ space. If the
    frontier has a non-convex kink — two Pareto-optimal labels $(t_1, c_1)$
    and $(t_2, c_2)$ where no scalar $\lambda$ makes one strictly preferred
    over the other — then \emph{no single $\lambda$-scalarised Dijkstra run
    will discover both labels}.

    This is precisely why the multi-label label-correcting algorithm is
    necessary. Repeated scalarised Dijkstra (varying $\lambda$ over a grid)
    would miss non-convex Pareto-optimal paths. The multi-label approach
    maintains the full frontier at each node simultaneously, recovering all
    Pareto-optimal labels regardless of the frontier's geometry.

    \textbf{Concrete example.} Three paths from $s$ to $t$:
    $(60\text{min}, 100\text{ crowd})$,
    $(75\text{min}, 40\text{ crowd})$,
    $(90\text{min}, 35\text{ crowd})$.
    The middle path is Pareto-dominated only if
    $75\text{min} - 60\text{min} = 15\text{min}$ of extra time is
    ``worth'' saving $60$ crowd units, while the third path costs 15 more
    minutes to save only 5 more crowd units. For intermediate $\lambda$,
    the second path lies below the line segment from first to third —
    a non-convex kink that no $\lambda$ scalarisation selects.
\end{remark}

\subsection{Critical Limitations}

\begin{remark}[Distributional Asymmetry]
    Transit delay distributions are \textbf{right-skewed}: they are bounded
    below by zero (a train cannot arrive before schedule beyond small margins)
    but are unbounded above due to service disruptions, signal failures, and
    cascade delays. Under right skew, variance as a risk measure is
    \textbf{economically inappropriate}: it penalises early arrivals
    symmetrically with late arrivals, despite the fact that arriving early
    incurs zero welfare cost to the passenger.
\end{remark}

\begin{remark}[Superior Risk Proxies]
    Two economically superior alternatives are:
    \begin{enumerate}
        \item \textbf{Semi-variance} (downside variance):
              \[
                  \mathrm{SV}[\crowd] = \E\left[\bigl(\crowd - \E[\crowd]\bigr)_+^2\right],
              \]
              where $(x)_+ = \max(x, 0)$. This penalises only crowd outcomes
              exceeding the mean, correctly reflecting passenger disutility.

        \item \textbf{Conditional Value at Risk (CVaR)} at confidence level
              $\alpha$ (e.g., $\alpha = 0.95$):
              \[
                  \CVaR_\alpha[\crowd] = \E\left[\crowd \mid \crowd \ge \mathrm{VaR}_\alpha[\crowd]\right].
              \]
              CVaR is coherent in the sense of Artzner et al.~\cite{artzner1999}
              (satisfying subadditivity, monotonicity, translation invariance,
              and positive homogeneity), whereas variance is not coherent.
              CVaR directly targets the worst-$1{-}\alpha$ fraction of crowding
              outcomes, which aligns with the scheduling objective of avoiding
              crush-load conditions at the tail.
    \end{enumerate}
\end{remark}

\begin{remark}[Practical Implication for the Engine]
    The current implementation uses $\lambda \cdot \crowd$ as a linear proxy.
    A more rigorous extension would replace this with the empirical CVaR
    computed over historical AFC gate count distributions for each
    (station, time-of-day, day-of-week) cell. This is left as future work
    and noted here to demonstrate awareness of the Markowitz analogy's limits
    to any reviewing quantitative researcher.
\end{remark}

% ══════════════════════════════════════════════════════════════════════════
\section{Why Label-Correcting is Necessary: Counterexample}
% ══════════════════════════════════════════════════════════════════════════

\begin{proposition}[Label-Setting Fails for Bi-Criteria Routing]
    \label{prop:label-setting-fails}
    In the bi-criteria (time, crowd) setting, a label-setting algorithm that
    \emph{permanently settles} a node on first extraction from the priority
    queue will miss Pareto-optimal labels in general.
\end{proposition}

\begin{proof}[Proof by counterexample]
    Consider two nodes $u$ and $v$ with two directed edges:
    \[
        e_1: \; u \to v, \quad t_{\text{dep}}=10, \; \tau=10, \; \crowd=100
        \quad \Rightarrow \quad (t_{\text{arr}}, \text{crowd}) = (20, 100)
    \]
    \[
        e_2: \; u \to v, \quad t_{\text{dep}}=20, \; \tau=10, \; \crowd=1
        \quad \Rightarrow \quad (t_{\text{arr}}, \text{crowd}) = (30, 1)
    \]
    Both labels $(20, 100)$ and $(30, 1)$ are Pareto-optimal: the first
    arrives earlier with higher crowd, the second arrives later with far
    lower crowd. Neither dominates the other.

    A label-setting algorithm extracts $(20, 100)$ first (smaller arrival
    time) and marks node $v$ as \emph{permanently settled}. It will never
    revisit $v$ to record $(30, 1)$, despite $(30, 1)$ being Pareto-optimal
    for any user with $\lambda > \frac{10}{99} \approx 0.101$
    (i.e., anyone willing to wait 10 minutes to avoid crowd cost of 99 units).

    The bi-criteria Pareto frontier at $v$ has size 2; label-setting returns
    size 1, missing the low-crowd Pareto-optimal path entirely.
\end{proof}

\begin{remark}[Why Standard Single-Criteria Dijkstra Does Not Have This Problem]
    In the single-criteria case, $f: V \to \N$ assigns exactly one optimal
    cost per node. Once a node is settled with the minimum-cost label, no
    future relaxation can improve it (by the non-negative edge weight property
    and FIFO). In the bi-criteria case, ``minimum cost'' is not a total order:
    a settled node may have multiple non-dominated labels, and a later
    extraction from the queue may carry a genuinely new Pareto-optimal label
    with worse time but better crowd.

    This is why the implementation is \textbf{label-correcting}: nodes may be
    re-inserted into the priority queue multiple times (up to $k$ times each,
    once per Pareto-optimal label), and the lazy deletion filter is used to
    skip stale extractions rather than permanently closing nodes.
\end{remark}

% ══════════════════════════════════════════════════════════════════════════
\section{Bounded-Wait Lookahead: Optimality Conditions}
% ══════════════════════════════════════════════════════════════════════════

\begin{proposition}[Bounded-Wait Suboptimality Conditions]
    \label{prop:lookahead-suboptimality}
    The Bounded-Wait Lookahead policy (select the minimum-composite departure
    within $[t, t + W_{\max}]$ among the next $k$ departures to destination $v$)
    may miss Pareto-optimal paths under two conditions:

    \begin{enumerate}
        \item \textbf{Window too tight}: The true Pareto-optimal departure for
              some edge $(u,v)$ on a Pareto-optimal path has
              $t_{\text{dep}} > t + W_{\max}$. Since the algorithm never waits
              beyond $W_{\max}$, this path is permanently unreachable regardless
              of $\lambda$.

        \item \textbf{Composite minimisation misses connection}: A departure
              $e^*$ that minimises composite at $(u,v)$ arrives at $v$ late
              enough to miss the only onward connection $(v,w)$ within the
              $W_{\max}$ window at $v$. A different departure $e'$ with higher
              composite at $(u,v)$ would have connected, producing a Pareto-
              optimal end-to-end path that $e^*$ cannot achieve.
    \end{enumerate}
\end{proposition}

\begin{remark}[Sufficiency for Optimality]
    The algorithm is globally Pareto-optimal if and only if for every edge
    $(u,v)$ on every Pareto-optimal path, the optimal departure falls within
    $[t, t + W_{\max}]$ and is among the $k$ departures enumerated. For
    BMRCL with 6-minute frequency and $W_{\max} = 1800\text{s}$ (30 min),
    the window covers approximately 5 departures, so $k=5$ is sufficient
    provided the network does not exhibit worst-case cascade misses.

    In practice, the synthetic pen $= 0$ in the scaffold means composite
    = $\tau + \lambda \cdot \crowd$ is time-independent, and no cascade miss
    can occur — the algorithm is optimal for the scaffold. The suboptimality
    risk is present only when non-zero time-dependent penalties are introduced.
\end{remark}

% ══════════════════════════════════════════════════════════════════════════
\section{Memory Architecture: Zero-Allocation Routing Proof}
% ══════════════════════════════════════════════════════════════════════════

\begin{proposition}[Allocation-Free Label Routing Invariant]
    After the first warm-up query, all subsequent calls to
    \texttt{ParetoDijkstra::run} make zero calls to the OS heap allocator
    for \texttt{Label} objects. One small allocation per call is made for the
    \texttt{QueryResult} vector shell (not for any \texttt{Label}).
\end{proposition}

\begin{proof}[Proof sketch]
    \begin{enumerate}
        \item \textbf{Label allocation.} All \texttt{Label} objects are
              allocated from \texttt{ArenaAllocator}, a fixed-capacity
              pre-allocated slab. \texttt{allocate()} pops from an intrusive
              free list or advances a bump pointer — $O(1)$, no system call.
              \texttt{reset()} restores the free list in $O(1)$.

        \item \textbf{Priority queue.} \texttt{pq\_} is a private member
              cleared via \texttt{while(!pq\_.empty()) pq\_.pop()} which reduces
              size but never deallocates the internal vector. After the first
              warm-up call reaches peak label count, the vector capacity is
              never exceeded and no reallocation occurs.

        \item \textbf{Pareto sets.} \texttt{result\_pareto\_} is a private
              member \texttt{vector<ParetoLabelSet>} sized to \texttt{num\_nodes}
              at construction. Each \texttt{ParetoLabelSet::labels\_} grows
              during the first warm-up call to its peak size ($\le k_{\max}$)
              and is never re-allocated thereafter.

        \item \textbf{Working buffers.} \texttt{best\_time\_} (reset via
              \texttt{std::fill}) and \texttt{dest\_scratch\_} (cleared via
              \texttt{clear()}) are private members pre-allocated at
              construction. No per-call heap activity.

        \item \textbf{QueryResult return (one allocation per call).}
              \texttt{result\_pareto\_} is moved into \texttt{QueryResult}
              at the end of \texttt{run()}. \texttt{std::move} leaves the
              source vector with \texttt{capacity() == 0} on all standard
              library implementations; the subsequent \texttt{resize()} must
              allocate a fresh buffer of $\sim|V| \times 24$ bytes
              ($\approx 2.4$\,KB for BMRCL). This is \emph{one allocation
              per call for the vector shell only} — no \texttt{Label} objects
              are heap-allocated. This does not affect the correctness or the
              latency of the routing algorithm itself; the allocation occurs
              at function exit, after the RDTSCP measurement closes.

              \begin{note}
              To eliminate this allocation entirely, the API could be changed
              to accept a \texttt{QueryResult\&} output parameter, allowing
              \texttt{result\_pareto\_} to be swapped into the caller's buffer
              in $O(1)$ with zero allocation. This is left as a future
              improvement.
              \end{note}
    \end{enumerate}

    The \textbf{Label allocation invariant} therefore holds strictly after
    warm-up: zero \texttt{malloc}/\texttt{new} calls for any \texttt{Label}
    object in any timed sample. The single per-call vector-shell allocation
    is an accepted cost of the by-value \texttt{QueryResult} API and does
    not appear in the RDTSC-measured routing latency.
\end{proof}

% ══════════════════════════════════════════════════════════════════════════
\section{Complexity Analysis}
% ══════════════════════════════════════════════════════════════════════════

\begin{proposition}[Time Complexity]
    \label{prop:complexity}
    The multi-label correcting Pareto-Dijkstra algorithm has time complexity
    \[
        O\!\left(k \cdot |E| \cdot \bigl(k + \log(k \cdot |V|)\bigr)\right),
    \]
    where $k$ is the maximum Pareto frontier size at any single node.
    The $k^2\cdot|E|$ term comes from \texttt{insert\_and\_dominate} forward
    pruning; the $k\cdot|E|\cdot\log(k\cdot|V|)$ term from heap operations.

    \begin{note}
    The previously stated bound $O(k\cdot(|E|+|V|)\log|V|)$ is incorrect for
    the label-correcting setting. Each of the $k$ Pareto labels per node can be
    re-inserted into the priority queue on every edge relaxation. The heap holds
    at most $O(k \cdot |V|)$ entries \emph{simultaneously} (at most $k$ labels
    per node, $|V|$ nodes); the log factor is therefore $O(\log(k\cdot|V|))$.
    Total heap insertions over the run are $O(k \cdot |E|)$, giving heap cost
    $O(k\cdot|E|\cdot\log(k\cdot|V|))$.

    Additionally, \texttt{insert\_and\_dominate}'s forward pruning (Step 3)
    costs $O(k)$ per call and is invoked $O(k\cdot|E|)$ times, contributing
    $O(k^2\cdot|E|)$ pruning work. The \emph{full} complexity is:
    \[
        O\!\left(k \cdot |E| \cdot \bigl(k + \log(k \cdot |V|)\bigr)\right).
    \]
    For BMRCL ($k\le 16$, $|E|\approx 300$, $|V|\approx 60$):
    pruning $\approx 76{,}800$ ops, heap $\approx 62{,}400$ ops.
    The pruning term dominates; quoting only the heap term understates the cost.

    The $O((|E|+|V|)\log|V|)$ factor is correct only for single-label Dijkstra
    where each node is settled exactly once; in the label-correcting case,
    nodes are revisited up to $k$ times.
    \end{note}

    For Namma Metro ($|V| \approx 60$, $|E| \approx 300$, $k \le 16$
    empirically from the Pareto frontier size bound in \Cref{prop:frontier-bound}),
    this yields approximately $300 \times 16 \times \log(16 \times 60)
    \approx 300 \times 16 \times 10 = 48{,}000$ heap operations per query —
    consistent with sub-millisecond latency targets at GHz clock rates.
\end{proposition}

\begin{proposition}[Pareto Frontier Size Bound]
    \label{prop:frontier-bound}
    At any node $v$, the Pareto frontier $\mathcal{L}(v)$ has at most
    $k \le \min(M, T)$ labels, where $M$ is the number of distinct
    crowd-cost levels and $T$ is the number of distinct arrival-time values
    reachable at $v$.

    For BMRCL with 5 daily service windows and 3 crowd tiers, $k \le 15$
    in practice. The worst-case graph structure is a \emph{fan graph}: $k$
    parallel paths from $s$ to $v$, each with strictly different
    (arrival-time, crowd-cost) pairs and no pair dominating another. This
    yields $|\mathcal{L}(v)| = k$ and is achieved when all paths are
    Pareto-incomparable.

    \textbf{Arena sizing rule:} The arena capacity must satisfy
    $C \ge |V| \times k_{\max}$, where $k_{\max}$ is the maximum frontier
    size. For BMRCL: $60 \times 16 = 960 \ll 65{,}536$. For Indian Railways
    ($|V| \approx 8{,}000$, $k_{\max} \approx 16$): $128{,}000 > 65{,}536$
    --- the arena would need resizing for nationwide deployment.
\end{proposition}

% ── Bibliography ──────────────────────────────────────────────────────────
\bibliographystyle{plain}
\begin{thebibliography}{9}

\bibitem{markowitz1952}
H.~Markowitz,
``Portfolio Selection,''
\textit{Journal of Finance}, vol.~7, no.~1, pp.~77--91, 1952.

\bibitem{artzner1999}
P.~Artzner, F.~Delbaen, J.-M.~Eber, and D.~Heath,
``Coherent Measures of Risk,''
\textit{Mathematical Finance}, vol.~9, no.~3, pp.~203--228, 1999.

\bibitem{delling2009}
D.~Delling, P.~Sanders, D.~Schultes, and D.~Wagner,
``Engineering Route Planning Algorithms,''
in \textit{Algorithmics of Large and Complex Networks}, Springer, 2009.

\bibitem{intel2010}
Intel Corporation,
``How to Benchmark Code Execution Times on Intel\textsuperscript{\textregistered}
IA-32 and IA-64 Instruction Set Architectures,''
White Paper, 2010.

\end{thebibliography}

\end{document}
